\documentclass[12pt,a4paper]{article}
\usepackage[margin=1in]{geometry}
\usepackage{amsmath,amssymb,amsthm}
\usepackage{hyperref}
\usepackage{enumitem}

\newtheorem{theorem}{Theorem}[section]
\newtheorem{corollary}[theorem]{Corollary}
\newtheorem{lemma}[theorem]{Lemma}
\newtheorem{definition}[theorem]{Definition}
\newtheorem*{remark}{Remark}

\title{\textbf{The Cascade Series --- Part III}\\[6pt]
General Relativity, Four Dimensions,\\
and Lorentzian Signature from the Cascade:\\
Lovelock Uniqueness, Spinor Compatibility,\\
and the Cosmological Constant}
\author{RTAC}
\date{March 2026}

\begin{document}
\maketitle

\begin{abstract}
The sphere-area cascade [1] derives a geometric invariant $I = 1.0996 \times 10^{-120}$
from orthogonality alone. The companion paper [2] shows that the cascade geometry
produces complex quantum mechanics via a forced precession angle $\alpha = \pi/2$. This
paper derives, from the same orthogonality axiom and with no additional empirical
input: (1) general relativity with cosmological constant $\Lambda = I$; (2) the observer's
spacetime dimension $d = 4$; and (3) Lorentzian metric signature $(-,+,+,+)$.

The argument has four components. \textbf{The cascade metric} (Section 4): the slicing
recurrence produces a foliated FRW-type metric with a natural lapse function and
extrinsic curvature. \textbf{Lovelock uniqueness} (Section 3): the Einstein equation with
cosmological constant is the unique divergence-free symmetric rank-2 metric equation
in exactly four spacetime dimensions. \textbf{Dimension derivation} (Section 9): $d = 4$ is the
unique dimension satisfying simultaneously (C1) the Clifford algebra classification
requires irreducibly complex spinors (forced by the cascade's $J^2 = -\text{Id}$), and (C2)
Lovelock uniqueness holds. The intersection of these conditions is $\{4\}$. A third
independent characterisation is established in Corollary 9.4: the Lorentzian cascade
metric has Ricci scalar $R^{(n)} = (n-1)(n-4)/a^4$, vanishing uniquely at $n = 4$.
\textbf{Lorentzian signature} (Section 10): the cascade produces both a Euclidean geometry
and a quantum propagator $e^{-i\lambda x}$ with $\lambda > 0$ from the forced precession. Given the
physical identification hypothesis, the propagator $e^{-iHt}$ with $H > 0$ is oscillatory
and unitary: Lorentzian time evolution. The signature $(-,+,+,+)$ follows without
semiclassics or analytic continuation.

The cosmological constant $\Lambda = I$ is derived, not fitted. The cascade's dark
energy equation of state is $w = -1$ exactly: $\Lambda = I$ is a fixed geometric constant
with no time-evolution mechanism, so $w = -1$ by definition. The Gauss--Bonnet
correction at $d = 5$ vanishes identically by two independent mechanisms (totally
umbilic cancellation and parity of the residual integrand), confirming the prediction is
structurally stable. The supplement (Section 14) proves the compactification radius
at each step is $R_{\text{eff}}(d) = 1/\sqrt{d + 3}$ exactly (from the Beta function), and presents the
exact Lorentzian Gauss--Codazzi analysis: the cascade metric contributes $w = 1/3$
(radiation) in cascade time, and $K_{\mu\nu} = 0$ at every equatorial embedding eliminates
the extrinsic curvature mechanism. These results resolve the question of the cascade's
matter fraction: the geometric route via Gauss--Codazzi is definitively closed, and
the matter fraction $\Omega_m$ is a topological quantity determined by the Bott partition of
cascade layers, not by the cascade metric geometry.
\end{abstract}

\tableofcontents
\newpage

\section{Division of Labour and Summary of the Series}

\subsection{The architecture}

The three-paper core rests on one hypothesis and three uniqueness theorems.

The hypothesis (this paper, Section 2): the cascade's abstract geometry is our physics.
Without it, the series is pure mathematics; with it, every physical result is a prediction
that tests the hypothesis.

The structure: the cascade provides geometric content; classical uniqueness theorems
identify the unique physical theory consistent with that content.

\begin{center}
\begin{tabular}{lll}
\hline
Paper & Cascade provides & Uniqueness provides \\
\hline
I [1] & $\Omega_d$ cascade, $I \approx 10^{-120}$ & --- (pure geometry) \\
II [2] & Complex state space, propagator $e^{-iHt}$ & Sphere geometry (Born rule) \\
III (this) & Metric, $\Lambda = I$, $d = 4$, $(-,+,+,+)$ & Lovelock, Clifford, propagator \\
\hline
\end{tabular}
\end{center}

\subsection{What this paper proves}

Given the physical identification hypothesis (Section 2), this paper derives four results
from the orthogonality axiom of [1]:

\begin{enumerate}
\item The cascade produces a foliated spacetime geometry with a natural lapse function
$N(d) = \sqrt{\pi} \cdot R(d)$ and a FRW-type 4D metric.
\item By Lovelock's theorem, the unique gravitational equation available to a 4D observer
is the Einstein equation; the cosmological constant equals the cascade invariant
$\Lambda = I$.
\item $d = 4$ is the unique spacetime dimension satisfying both the cascade's forced complex
spinor requirement and Lovelock uniqueness.
\item Lorentzian metric signature $(-,+,+,+)$ is the unique signature consistent with the
cascade's forced propagator $e^{-iHt}$ with $H > 0$.
\end{enumerate}

\section{The Physical Identification Hypothesis}

The cascade is a mathematical object: a sequence of sphere areas $\{\Omega_d\}$ derived from
orthogonality. Paper I [1] is pure mathematics; it makes no physical claim beyond recording
the numerical coincidence $I \approx \rho_\Lambda/M_{\text{Pl}}^4$. Paper II [2] shows the cascade's state space has the
structure of quantum mechanics. This paper derives the cascade's gravitational structure.
Both Papers II and III require one physical input.

\begin{definition}[Physical identification hypothesis]
The infinite-dimensional unit ball,
descended to four dimensions, is indistinguishable from our universe. Concretely: the
cascade metric is the spacetime metric; the cascade propagator is the physical time-evolution
operator; the cascade invariant $I$ is the cosmological constant. This hypothesis is the
claim that the cascade's abstract geometry is our physics. Without it, the series is pure
mathematics; with it, every physical result is a prediction that tests the hypothesis.
\end{definition}

\begin{remark}[Uniformity of the physical identification]
Definition 2.1 maps cascade
content to physical energy uniformly across all layer types. By Corollary 4.8a of [1], sphere
area $\Omega_{d-1}$ is the only independent cascade quantity at level $d$; every other cascade object
is derived from sphere areas. Therefore the identification assigns energy $C \cdot \Omega_{d-1}$ at each
level with $C$ a universal constant---no layer-type-dependent modulation is available, because
no independent geometric quantity exists to carry such modulation. Matter and vacuum
are observer labels applied to the Bott partition of cascade layers: the propagator phase
classification ([2], Corollary 6.5; refined to period 8 via Clifford in this paper, Section 9) establishes the partition topologically, and the fermion
generation structure ([3]) identifies which sector is matter via the hairy ball zeros and
observed generations. The cascade geometry itself makes no distinction between the sectors.
\end{remark}

\begin{remark}
The physical identification hypothesis does not assume: the dimension of
spacetime (derived as $d = 4$ in Section 9); the metric signature (derived in Section 10);
the form of the gravitational equation (derived via Lovelock); or the value of $\Lambda$ (derived as
$I$ from [1]). It assumes only that the cascade corresponds to something physical. This is
the minimal assumption needed to connect abstract geometry to physics.
\end{remark}

\section{Lovelock's Theorem and the Privilege of Four Dimensions}

\begin{theorem}[Lovelock {\cite{lovelock1,lovelock2}}]
Let $(M, g)$ be a smooth $d$-dimensional Lorentzian manifold
with $d = 4$. The most general symmetric divergence-free rank-2 tensor constructed from
$g_{\mu\nu}$ and its first and second derivatives is
\[
E^{\mu\nu} = a\,G^{\mu\nu} + b\,g^{\mu\nu},
\]
where $G^{\mu\nu}$ is the Einstein tensor and $a$, $b$ are constants. Setting $a = 1$, $b = -\Lambda$ recovers
the Einstein equation with cosmological constant.
\end{theorem}

The theorem is dimension-dependent. In $d$ spacetime dimensions, the most general
such tensor is a sum of $\lfloor d/2 \rfloor$ independent Lovelock terms plus $\Lambda$.

\begin{center}
\begin{tabular}{lcl}
\hline
Dimension $d$ & Lovelock terms & Gravity unique? \\
\hline
3 & $\Lambda$ only & No propagating gravity \\
4 & $G^{\mu\nu} + \Lambda g^{\mu\nu}$ only & Yes: GR is unique \\
5, 6 & Einstein + Gauss--Bonnet + $\Lambda$ & No: free coupling \\
$\geq 7$ & Further Lovelock terms + $\Lambda$ & No \\
\hline
\end{tabular}
\end{center}

\begin{remark}
Lovelock uniqueness at $d = 4$ is a theorem about tensor equations, not a
consequence of the cascade. The cascade's role is to provide the metric that Lovelock needs
as input. Section 9 derives $d = 4$ from the consistency of the cascade's quantum structure
with Lovelock uniqueness.
\end{remark}

\section{The Cascade Metric}

\subsection{Foliation from slicing}

The cascade's slicing recurrence decomposes $B^{d+1}$ into $d$-balls parametrised by the
perpendicular coordinate $x \in [-1, 1]$. Each cross-section at height $x$ is a $d$-ball of radius
$r(x) = \sqrt{1 - x^2}$. This is a foliation; $x$ plays the role of time (identified in [2], Section 7 via
the irreversibility of slicing). The full $(d + 1)$-dimensional metric is:
\[
ds^2 = dx^2 + (1 - x^2)\,d\sigma_{d-1}^2.
\]

\subsection{The lapse function}

The ratio of adjacent volumes is the cascade's lapse:
\[
N(d) = \sqrt{\pi} \cdot R(d), \qquad R(d) = \frac{\Gamma((d + 1)/2)}{\Gamma((d + 2)/2)} \approx \sqrt{\frac{2\pi}{d}} \quad (d \gg 1).
\]
At $d = 4$: $N(4) = \sqrt{\pi} \cdot \Gamma(5/2)/\Gamma(3) = 3\pi/8 \approx 1.178$. As $d \to \infty$: $N \to 0$ (no temporal
separation at the cascade's geometric origin).

\subsection{The extrinsic curvature}

The extrinsic curvature of the leaf at height $x$ is $K = -x/\sqrt{1 - x^2}$. At the equator $x = 0$:
$K = 0$. This is the cascade's Hubble parameter.

\subsection{The 4D projected metric}

Projecting onto a 4-dimensional subspace, the 4D observer's metric is:
\[
ds^2_{\text{4D}} = -dt^2 + a(t)^2\,d\sigma_3^2.
\]
This is the FRW metric with unit lapse. Isotropy follows from the cascade's spherical
symmetry; homogeneity from the uniform slicing recurrence. The Lorentzian sign $g_{tt} = -1$
is derived in Section 10.

\section{The Cosmological Constant}

\subsection{From cascade invariant to $\Lambda$}

By [1], Theorem 9.2:
\[
I = \Omega_{19} \times \Omega_{217} \times \frac{198}{217} = 1.0996 \times 10^{-120}.
\]
The 4D observer's vacuum energy satisfies $\rho_\Lambda/M_{\text{Pl}}^4 \approx 1.1 \times 10^{-120}$. The identification $\Lambda = I$
is not a fit: $I$ is fixed by $\pi$ independently of any physical measurement. The agreement is
$0.04\%$ under the reduced-Planck-mass convention.

\subsection{Why $\Lambda$ is not a free parameter}

In standard GR, $\Lambda$ is a free parameter: Lovelock permits any $b$ in $E^{\mu\nu} = aG^{\mu\nu} + bg^{\mu\nu}$. The
cascade removes this freedom. The constant $b = -\Lambda$ is fixed by the cascade invariant $I$; the
observer does not choose it. The cosmological constant problem---why $\Lambda \sim 10^{-120}M_{\text{Pl}}^4$---is
answered by the hierarchy theorem of [1]: $\log_{10}(1/\Omega_{d_2}) \approx \pi e^{2\sqrt{\pi}}(2\sqrt{\pi} - 1)/\ln 10 \approx 120.3$
orders of magnitude arise from the Gamma function's superexponential growth between
$d_1 = 19$ and $d_2 = 217$.

\section{The Dark Energy Equation of State}

\begin{theorem}[Dark energy equation of state]
The cascade predicts $w = -1$ exactly.
\end{theorem}

\begin{proof}
$\Lambda = I$ is a fixed geometric constant, determined entirely by $\pi$ through the cascade's
threshold structure ([1], Theorem 9.2). A fixed cosmological constant has $\rho_\Lambda = \text{const}$,
giving $p_\Lambda = -\rho_\Lambda$ and $w = p/\rho = -1$ by definition. No time-evolution mechanism exists
in the cascade for $\Lambda$: the thresholds $d_1 = 19$ and $d_2 = 217$ are permanent features of the
Gamma function, not snapshots of a dynamical field. The lapse function $N(4) = 3\pi/8$ is a
fixed Gamma-function value, not a cosmological-time-evolving quantity.
\end{proof}

\subsection{Gauss--Bonnet stability}

The cascade geometry at $d = 5$ has $S^3$ cross-sections that are totally umbilic: $K_{ij} =
-(x/\sqrt{1 - x^2})\,h_{ij}$ exactly, so $K_{ij} = (K/3)h_{ij}$. For totally umbilic hypersurfaces, the leading
Gauss--Bonnet boundary terms cancel identically:
\[
2K_{ij}K^{ij}K - \tfrac{2}{3}K^3 = \tfrac{2}{3}K^3 - \tfrac{2}{3}K^3 = 0.
\]
The residual GB term $2K\tilde{R} - 2K_{ij}\tilde{R}^{ij} \propto -24x/(1 - x^2)^2$ is odd in $x$; the cascade measure
$(1 - x^2)^{3/2}$ is even. The integral vanishes exactly:
\[
\int_{-1}^{1} (1 - x^2)^{3/2} \cdot \frac{-24x}{(1 - x^2)^2}\,dx = -24\int_{-1}^{1} \frac{x}{\sqrt{1 - x^2}}\,dx = 0.
\]
The bulk GB contribution is suppressed by $H_0^2/M_{\text{Pl}}^2 \sim 10^{-120}$.

\begin{corollary}
$w = -1$ receives zero Gauss--Bonnet correction at leading order, by two
independent mechanisms: (i) totally umbilic cancellation from the cascade's spherical
symmetry; (ii) parity of the residual integrand under symmetric slicing. Both follow from
the orthogonality axiom. The prediction is structurally protected, not merely numerically
small.
\end{corollary}

\begin{remark}[Quantitative cosmology]
The cascade's geometric parameters determine the
background cosmological observables. The density fractions $\Omega_\Lambda = (\pi-1)/\pi$,
$\Omega_m = 1/\pi$ (leading order) or $\Omega_m^{\rm Bott} = 0.31150$ (subleading), and
$\Omega_b = 1/(2\pi^2)$ follow from the cascade's sphere-area structure. The Hubble constant
$H_0 = 71.05$~km/s/Mpc follows from the lapse-corrected Friedmann equation with
$N(4) = 3\pi/8$, sitting between the Planck measurement (67.4) and the SH0ES value (73.0).
These parameters predict a sound horizon $r_d \approx 141$~Mpc. The apparent DESI preference
for $w \neq -1$ is a zero-parameter consequence of imposing Planck's sound horizon
$r_d = 147.6$~Mpc on a universe where $r_d \approx 141$~Mpc: forcing the wrong ruler onto
cascade distances produces apparent $w \approx -0.80$, matching DESI's BAO+CMB-only report
of $-0.76 \pm 0.11$ within $1\sigma$. With the cascade's ruler, DESI BAO is consistent with
$w = -1$ throughout.
\end{remark}

\begin{theorem}[Orthogonality of coupling and vacuum sectors]
The cascade lapse factorises
as $N(d) = \sqrt{\pi} \cdot R(d)$. Gauge couplings satisfy $\alpha(d) = N(d)^2/(4\pi) = R(d)^2/4$: the
factor of $\pi$ from squaring $\sqrt{\pi}$ cancels the denominator exactly, leaving $\alpha = R(d)^2/4$: the
multiplicative $\sqrt{\pi}$ factor of the lapse drops out exactly. The cosmological constant $\Lambda = I$ is
generated by the threshold conditions $p(d) = c_1 = \tfrac{1}{2}\ln\pi$ and $p(d) = c_2 = \sqrt{\pi}$, which are
determined by the constant part of $p(d)$---the same $\sqrt{\pi}$ factor---with no dependence on $R(d)$
at the threshold dimensions. The two sectors therefore draw on orthogonal factors of $N(d)$:
the multiplicative $\sqrt{\pi}$ of the lapse cancels from $\alpha$ and generates $\Lambda$; $R(d)$ determines $\alpha$ and
does not enter the threshold conditions. This is why $\Lambda$ does not run while couplings run.
\end{theorem}

\begin{proof}
$\alpha(d) = N(d)^2/(4\pi) = \pi R(d)^2/(4\pi) = R(d)^2/4$. The $\sqrt{\pi}$ factor cancels entirely;
$\alpha$ depends only on $R(d)$. The thresholds $d_1 = 19$ and $d_2 = 217$ are determined by
$p(d) = c_1 = \tfrac{1}{2}\ln\pi$ and $p(d) = c_2 = \sqrt{\pi}$, where $c_1$, $c_2$ are generated from $\sqrt{\pi}$ alone ([1],
Theorem 6.5). $\Lambda = I$ is a function of $\Omega_{d_1}$ and $\Omega_{d_2}$, which are values of the sphere-area
sequence at the threshold dimensions; these values are determined by $\sqrt{\pi}$ through the
threshold conditions. At those dimensions, $R(d_1)$ and $R(d_2)$ enter only the $d$-dependent
decay rate, not the threshold conditions themselves. Therefore $\Lambda$ contains no dependence
on $R(d)$ through the threshold structure.
\end{proof}

\section{Matter from Extra Dimensions}

\subsection{The Kaluza--Klein perspective}

The cascade produces a 217-dimensional geometry. The 4D observer has integrated out
213 directions via the slicing recurrence. The purpose of this section is to identify the
effective stress-energy tensor $T^{\mu\nu}_{\text{eff}}$ that enters the Einstein equation projection identity
(Section 8): Lovelock's theorem forces $G^{\mu\nu} + \Lambda g^{\mu\nu} = 8\pi G\,T^{\mu\nu}$, and $T^{\mu\nu}_{\text{eff}}$ is the cascade's
geometric account of what fills the right-hand side.

\begin{remark}[Induced matter via Gauss--Codazzi]
Let $G^{\mu\nu}_{(217)}$ be the Einstein tensor of
the cascade's 217-dimensional geometry and $G^{\mu\nu}_{(4)}$ be the 4D projected Einstein tensor. The
difference $8\pi G\,T^{\mu\nu}_{\text{eff}} = G^{\mu\nu}_{(4)} - [G^{\mu\nu}_{(217)}]_{\text{4D}}$ defines the effective stress-energy seen by the 4D
observer, following from the Gauss--Codazzi equations for the embedding $M^4 \subset M^{217}$. This
is a structural identification, not a dynamical calculation. Section 14 establishes exact
results for the Lorentzian cascade metric: $K_{\mu\nu} = 0$ at every equatorial embedding in the
chain $S^3 \hookrightarrow S^4 \hookrightarrow \cdots \hookrightarrow S^{216}$, and the single-step cascade geometry contributes $w = 1/3$
(radiation) in cascade time. The cascade's matter fraction $\Omega_m$ is determined by the Bott
partition of cascade layers, not by the cascade metric geometry; this is a topological rather
than geometric quantity.
\end{remark}

\section{Self-Consistency: Einstein Equation as Projection Identity}

The 4D observer has: a metric $g_{\mu\nu}$ (Section 4); $\Lambda = I$ (Section 5); $T^{\mu\nu}_{\text{eff}}$ from the extra
dimensions (Section 7); and Lovelock's theorem (Section 3). Lovelock then requires
$G^{\mu\nu} + \Lambda g^{\mu\nu} = 8\pi G\,T^{\mu\nu}$. This is a projection identity, not a postulate.

\begin{theorem}[Einstein equation as projection identity]
The Gauss--Codazzi equations
for $M^4 \subset M^{217}$, combined with Lovelock's theorem at $d = 4$, yield the Einstein equation
with $\Lambda = I$ and $T^{\mu\nu} = T^{\mu\nu}_{\text{eff}}$ as an identity. Every term is geometric output of the cascade;
no gravitational postulate is added.
\end{theorem}

\section{Deriving $d = 4$}

\begin{definition}[Cascade-compatible dimension]
A spacetime dimension $d$ is cascade-compatible if it satisfies:
\begin{itemize}
\item[(C1)] The minimal spinor representation of $\text{Spin}(1, d - 1)$ is irreducibly complex (not
realisable over $\mathbb{R}$).
\item[(C2)] The Einstein equation with cosmological constant is the unique divergence-free symmetric rank-2 tensor metric equation (Lovelock uniqueness, Theorem 3.1).
\end{itemize}
\end{definition}

\subsection{Why (C1) is forced by the cascade}

From Theorem 6.1 of [2]: the orthogonality axiom forces the precession angle $\alpha = \pi/2$,
producing the complex structure $J: e_1 \mapsto e_2$, $e_2 \mapsto -e_1$ with $J^2 = -\text{Id}$.

\begin{lemma}[Complex $J$ requires complex spinors]
For the cascade's state space (carrying
$J$ with $J^2 = -\text{Id}$) to serve as the spinor representation space of $\text{Spin}(1, d - 1)$, the minimal
spinor representation must be irreducibly complex. Majorana-real spinors are incompatible:
by Schur's lemma over $\mathbb{R}$, the commutant algebra of an irreducibly real representation is
$\mathbb{R}$, admitting no $J$ with $J^2 = -\text{Id}$.
\end{lemma}

\begin{proof}
If the spinor representation is Majorana (real), the representation space is $\mathbb{R}^n$ and
the commutant is $\mathbb{R}$ (Schur's lemma). Any $J$ commuting with all group elements satisfies
$J = c \cdot \text{Id}$ for some $c \in \mathbb{R}$. Then $J^2 = c^2\,\text{Id}$. Since $c^2 \geq 0$, we cannot have $J^2 = -\text{Id}$.
Contradiction.
\end{proof}

\subsection{The Clifford algebra classification}

The type of the minimal spinor of $\text{Spin}(1, d - 1)$ is determined by the Clifford algebra
$\text{Cl}(1, d - 1)$, following Bott periodicity with period 8.

\begin{center}
\begin{tabular}{lllll}
\hline
$d$ & $\text{Cl}(1, d - 1)$ structure & Min.\ spinor & Type & (C1) \\
\hline
2 & $M_2(\mathbb{R})$ & 2 real & Majorana & No \\
3 & $M_2(\mathbb{R}) \oplus M_2(\mathbb{R})$ & 2 real & Majorana & No \\
4 & $M_2(\mathbb{C}) \otimes_\mathbb{R} M_2(\mathbb{R})$ & 2 complex & Weyl & Yes \\
5 & $M_4(\mathbb{C})$ & 4 complex & Dirac & Yes \\
6 & $M_4(\mathbb{C}) \oplus M_4(\mathbb{C})$ & 4 complex & Weyl & Yes \\
7 & $M_8(\mathbb{R})$ & 8 real & Majorana & No \\
8 & $M_{16}(\mathbb{R})$ & 8 real & Majorana & No \\
\hline
\end{tabular}
\end{center}

Data from Lounesto [6], Lorentzian signature $(1, d - 1)$. Condition (C1) holds for
$d \in \{4, 5, 6\}$ in the first Bott period, then $\{10, 11, 12\}$, etc.

\subsection{The main theorem}

\begin{theorem}[Unique cascade-compatible dimension]
$d = 4$ is the unique dimension
satisfying both (C1) and (C2).
\end{theorem}

\begin{proof}
From the Clifford classification: (C1) holds for $d \in \{4, 5, 6, 10, 11, 12, \ldots\}$. From
Theorem 3.1: (C2) holds if and only if $d = 4$. Intersection: $\{4, 5, 6, 10, \ldots\} \cap \{4\} = \{4\}$.
\end{proof}

\begin{corollary}[Third characterisation of $d = 4$]
The Lorentzian cascade scale factor
$a(t) = \sqrt{1 - t^2}$ gives $\dot{a} = -t/a$ and $\ddot{a} = -1/a^3$. The Ricci scalar of the $n$-dimensional
FRW metric with $k = 1$ and this scale factor is:
\[
R^{(n)} = \frac{(n - 1)(n - 4)}{a^4}.
\]
This vanishes if and only if $n = 4$. The Lorentzian cascade metric is therefore Ricci-flat
uniquely in four spacetime dimensions, providing a third independent characterisation of
$d = 4$ complementing Lovelock uniqueness (C2) and complex spinor compatibility (C1).
\end{corollary}

\begin{proof}
For $n$-dimensional FRW with $k = 1$: $R_{tt} = -(n - 1)\ddot{a}/a = (n - 1)/a^4$ and
$R_{ij} = [(n - 2)(k + \dot{a}^2)/a^2 + \ddot{a}/a]\,g_{ij} = (n - 3)/a^4 \cdot g_{ij}$. The Ricci scalar: $R^{(n)} = g^{tt}R_{tt} + g^{ij}R_{ij} =
-(n - 1)/a^4 + (n - 1)(n - 3)/a^4 = (n - 1)(n - 4)/a^4$. At $n = 4$: $R^{(4)} = 3 \cdot 0/a^4 = 0$. For
$n \neq 4$: $(n - 1)(n - 4) \neq 0$.
\end{proof}

\begin{remark}
Without (C1): $d = 4$ is selected by Lovelock alone, but the quantum state
space may be incompatible with $J^2 = -\text{Id}$. Without (C2): dimensions 5, 6, 10,\ldots\ have
complex spinors but non-unique gravity with free coupling constants. Only $d = 4$ is
simultaneously quantum-compatible, gravitationally unique, and Ricci-flat under the cascade
scale factor.
\end{remark}

\begin{remark}[Division algebra interpretation]
The cascade forces the quantum amplitude
algebra to be $\mathbb{C}$ (real dimension 2). The associative normed division algebras over $\mathbb{R}$ are $\mathbb{R}$,
$\mathbb{C}$, $\mathbb{H}$, $\mathbb{O}$ (Hurwitz theorem [7]; $\mathbb{O}$ is non-associative and excluded from quantum mechanics
by associativity of sequential measurements). The minimal associative algebra over $\mathbb{C}$ is $\mathbb{H}$
(quaternions), of real dimension 4. The spacetime dimension equals the real dimension of
the next associative division algebra: $d = \dim_\mathbb{R}\,\mathbb{H} = 4$.
\end{remark}

\section{Lorentzian Signature from the Propagator}

The cascade geometry is Euclidean: $ds^2_E = dx^2 + (1 - x^2)\,d\sigma_{d-1}^2$ with positive-definite
metric. The cascade propagator is $K(x) = |K|\,e^{-i\lambda x}$ with $\lambda > 0$---already oscillatory for
real $x$. These are two descriptions of the same object under the physical identification
hypothesis. The Lorentzian signature is the unique metric sign consistent with both
descriptions simultaneously.

\subsection{The propagator characterises the signature}

\begin{lemma}[Spectral gap is positive]
The cascade spectral gap satisfies $\lambda_\infty =
\tfrac{1}{4}\psi^{(1)}(\sigma/2) > 0$ for all $\sigma > 0$, where $\psi^{(1)}$ is the trigamma function. This follows from $p'(d) =
\tfrac{1}{4}\psi^{(1)}((d + 1)/2) > 0$ (strict monotonicity of $p$).
\end{lemma}

From Theorems 6.1 and 7.1 of [2]: the forced precession $\alpha = \pi/2$ gives the cascade
propagator:
\[
K(t) = |K| \cdot e^{-iHt}, \qquad H = \lambda_\infty > 0,
\]
where $t = x$ is the physical time identified with the slicing coordinate.

\begin{theorem}[Lorentzian signature from orthogonality]
The spacetime metric of the
4D cascade observer has signature $(-,+,+,+)$.
\end{theorem}

\begin{proof}
\textbf{Temporal component.} Lorentzian signature means $g_{tt} < 0$. This is equivalent,
by definition, to time evolution being generated by a self-adjoint Hamiltonian with real
eigenvalues, giving unitary oscillatory evolution $e^{-iHt}$. The cascade propagator (from
Theorems 6.1 and 7.1 of [2]) is $K(t) = |K| \cdot e^{-iHt}$ with $H = \lambda_\infty > 0$, which is oscillatory
with $|K| = \text{const}$. Since the cascade propagator is oscillatory, the identification hypothesis
(Definition 2.1) gives $g_{tt} < 0$.

\textbf{Spatial components.} The cascade's state space is a complex Hilbert space $\mathcal{H}$ ([2]).
The Hilbert space inner product is positive-definite: $\langle\psi|\psi\rangle > 0$ for all $|\psi\rangle \neq 0$. The
inner product on equal-time surfaces requires a positive-definite spatial volume element
$\sqrt{\det g^{(3)}} > 0$, which holds if and only if $g_{ij} > 0$. The signature is $(-,+,+,+)$.
\end{proof}

\begin{remark}[Euclidean geometry, Lorentzian physics]
The cascade's underlying geometry is Euclidean; the physical spacetime is Lorentzian. These are not contradictory. The
connection is through the propagator's character: the cascade's Euclidean geometry produces
an oscillatory propagator (because of the forced precession), and oscillatory propagator
character is the physical content of Lorentzian signature. The Euclidean geometry generates
Lorentzian physics through the dynamics, not through analytic continuation of the metric.
\end{remark}

\begin{remark}[Three consequences of $\psi^{(1)} > 0$]
The strict positivity $\psi^{(1)}(x) > 0$ simultaneously establishes: (a) unique natural zero of $p(d)$ in [1]: requires $p'(d) > 0$; (b) positive
spectral gap $\lambda_\infty > 0$; (c) Lorentzian signature: $\lambda_\infty > 0$ makes the propagator oscillatory,
requiring $g_{tt} < 0$ (Theorem 10.2).
\end{remark}

\section{Scale Factors: $G$ and $\hbar$}

The cascade produces dimensionless quantities. Physical gravity requires $G$ (Newton) and
$\hbar$ (Planck). Both enter as unit-matching constants when geometric rates are identified
with physical observables:

\begin{itemize}
\item $G$: proportionality between 4D curvature and the effective stress-energy from the
extra dimensions. Set by the embedding geometry of $M^4 \subset M^{217}$.
\item $\hbar = (\pi/2)/\Delta t$: the ratio of the forced precession angle to the physical time increment
(Section 9 of [2]). The angle $\pi/2$ is derived; $\Delta t$ is not.
\end{itemize}

\section{Gravitons and Degrees of Freedom}

In $d$ spacetime dimensions, a massless graviton has $d(d - 3)/2$ physical polarisations. At
$d = 4$: 2 polarisations (the standard GR plus and cross modes). The remaining 23,245
polarisations of the 217-dimensional graviton are Kaluza--Klein modes: a tower of massive
particles appearing to the 4D observer as matter on adjacent shells. Their mass scale is
set by the compactification radii $R_{\text{eff}}(d) = 1/\sqrt{d + 3}$, derived in [1], Theorem 4.2.

\section{The Complete Derivation}

The full chain from the orthogonality axiom to the metric of general relativity:

\begin{center}
\small
\begin{tabular}{lll}
\hline
Step & Result & Source \\
\hline
Orthogonality $\Rightarrow \sqrt{\pi}$ & Cascade constant & [1], Thm 3.1 \\
$\sqrt{\pi} \Rightarrow I \approx 10^{-120}$ & Cascade invariant & [1], Thm 9.2 \\
Orthogonality $\Rightarrow J^2 = -\text{Id}$ & Complex structure & [2], Thm 6.4 \\
$J^2 = -\text{Id} \Rightarrow$ complex spinors (C1) & Schur + Clifford & Section 9 \\
(C1)$\cap$Lovelock (C2)$\Rightarrow d = 4$ & Dimension theorem & Thm 9.3 \\
$R^{(n)} = (n-1)(n-4)/a^4 \Rightarrow R^{(4)} = 0$ & Third char.\ of $d = 4$ & Cor 9.4 \\
$d = 4 \Rightarrow$ FRW metric + $\Lambda = I$ & Cascade geometry & \S4, \S5 \\
Physical identification hypothesis & Cascade geom = physics & Def 2.1 \\
Hypothesis + $\lambda_\infty > 0 \Rightarrow g_{tt} < 0$ & Propagator & Thm 10.2 \\
$d = 4$ + $(-,+,+,+) \Rightarrow$ Einstein eq. & Lovelock & Section 8 \\
Fixed $\Lambda = I \Rightarrow w = -1$ & Definition & Section 6 \\
GB corrections vanish & Two mechanisms & Section 6 \\
$N = \sqrt{\pi}\,R \Rightarrow$ orthogonal sectors & Couplings vs vacuum & Thm 6.4 \\
Extra dims $\Rightarrow T^{\mu\nu}$ & Gauss--Codazzi & Section 7 \\
\hline
\end{tabular}
\end{center}

\subsection{Full result table}

\begin{center}
\small
\begin{tabular}{lll}
\hline
Physical result & Cascade provides & Classical theorem \\
\hline
Complex QM, Born rule, $\hbar$ & Complex structure, propagator & Sphere geometry, [2] \\
Cosmological constant $\Lambda = I$ & Cascade invariant $I$ & Fixed, not free \\
$d = 4$ spacetime dimensions & Complex spinors + gravity & Clifford, Lovelock \\
$R^{(4)} = 0$: Ricci-flat cascade metric & Lorentzian scale factor & GR, Corollary 9.4 \\
Lorentzian signature $(-,+,+,+)$ & Propagator $e^{-iHt}$, $H > 0$ & Path integral identity \\
Einstein equation & Metric + $\Lambda = I$ & Lovelock \\
FRW cosmology & Cascade foliation & Geometry \\
$w = -1$ dark energy & Fixed $\Lambda = I$; GB vanishes & Section 6 \\
$\Lambda \approx 10^{-120}M_{\text{Pl}}^4$ & Sphere-area hierarchy & [1], Thm 10.1 \\
Orthogonal coupling and vacuum sectors & $N = \sqrt{\pi}\,R$ factorisation & Thm 6.4 \\
$K_{\mu\nu} = 0$: totally geodesic embedding & Lorentzian cascade metric & Section 14 \\
\hline
\end{tabular}
\end{center}

\section{Supplement: Asymptotic Compactification and the Cascade Metric}

This section integrates the compactification results of [1], Section 4 into the gravitational
framework, and presents the exact Lorentzian Gauss--Codazzi analysis.

\subsection{The compactification radius of each dimension}

\begin{theorem}[Compactification radius; {[1]}, Theorem 4.2]
At each slicing step $d$, the
integrated-out direction retains an effective radius
\[
R_{\text{eff}}(d) = \frac{1}{\sqrt{d + 3}},
\]
derived exactly from $\langle x^2 \rangle = B(3/2, d/2 + 1)/B(1/2, d/2 + 1) = 1/(d + 3)$.
\end{theorem}

The compactification is asymptotic: $R_{\text{eff}} \to 0$ as $d \to \infty$, but $R_{\text{eff}} \neq 0$ at any finite $d$.
There is no sharp boundary, no horizon, no topology change---only exponential suppression.

\begin{center}
\begin{tabular}{cccc}
\hline
$d$ & $R_{\text{eff}}$ & $N(d)$ & Physical role \\
\hline
4 & 0.37796 & 1.17810 & Observer dimension \\
5 & 0.35355 & 1.06667 & $S^3$ boundary layer \\
7 & 0.31623 & 0.91429 & Volume maximum $d_0$ \\
12 & 0.25820 & 0.70870 & SU(3) layer \\
13 & 0.25000 & 0.68198 & SU(2) breaking \\
19 & 0.21320 & 0.56755 & First threshold $d_1$ \\
217 & 0.06742 & 0.16997 & Second threshold $d_2$ \\
\hline
\end{tabular}
\end{center}

\begin{remark}[KK mass scale]
The KK mass scale at dimension $d$ is $M_{\text{KK}}(d) \sim
1/R_{\text{eff}}(d) = \sqrt{d + 3}$ in cascade units. No free parameter adjusts the mass scale.
\end{remark}

\subsection{Boundary dominance and the primacy of sphere areas}

\begin{theorem}[Boundary dominance; {[1]}, Theorem 4.4]
$\Omega_{d-1}/V_d = d$ for all $d \geq 1$.
\end{theorem}

At $d = 4$: $\Omega_3/V_4 = 4$, so the boundary $S^3$ carries $4/5 = 80\%$ of the content of $B^4$. At
$d = 217$: essentially all content is on the boundary. The cascade's content is its boundaries,
increasingly so at each step.

\subsection{The observer on $S^3$}

The 4D observer exists on the boundary $S^3$ of the $d = 5$ layer. Theorem 14.3 shows this
boundary carries $5/6 \approx 83\%$ of the $d = 5$ content. The observer's physics is a boundary
theory, not because holography is assumed, but because the boundary dominates the
volume at every cascade layer.

The $S^3$ horizon of a $d = 5$ Schwarzschild black hole has topology $S^3$, matching the 3
spatial dimensions. The cascade's $d = 5$ layer provides this horizon. The de Sitter horizon
area $A = 12\pi/\Lambda \sim 10^{120}$ in Planck units is the cascade hierarchy inverted: $\Omega_7/\Omega_{217} \approx 10^{121}$.

\begin{remark}[The cosmological constant as inverse boundary area]
With boundary
dominance, the identification $\Lambda = I$ acquires a geometric interpretation: $\Lambda$ measures the
inverse boundary area of the cascade's distinguished layers. The smallness of $\Lambda$ reflects the
largeness of the boundary---a high-capacity boundary ($S^3$ with $\Omega_3 = 2\pi^2 \approx 19.7$) encodes
an enormous cascade (213 compactified directions), giving $\Lambda = I \approx 10^{-120}$.
\end{remark}

\subsection{The Lorentzian Gauss--Codazzi analysis}

We install the Lorentzian cascade metric directly from Theorem 10.2, without analytic
continuation. The $n$-dimensional Lorentzian cascade metric is:
\[
g^L = -dt^2 + (1 - t^2)\,d\sigma_{n-1}^2, \qquad a(t) = \sqrt{1 - t^2}.
\]
With $\dot{a} = -t/a$ and $\ddot{a} = -1/a^3$, two exact identities hold:
\[
\frac{k + \dot{a}^2}{a^2} = \frac{1}{a^4}, \qquad \frac{\ddot{a}}{a} = -\frac{1}{a^4}.
\]

\begin{theorem}[Lorentzian cascade Ricci tensor]
For the $n$-dimensional Lorentzian cascade metric:
\[
R^{(n)}_{tt} = \frac{n - 1}{a^4}, \qquad R^{(n)}_{ij} = \frac{n - 3}{a^4}\,g_{ij}, \qquad R^{(n)} = \frac{(n - 1)(n - 4)}{a^4},
\]
\[
G^{(n)}_{tt} = \frac{(n - 1)(n - 2)}{2a^4}, \qquad G^{(n)}_{NN} = -\frac{(n - 2)(n - 5)}{2a^4}.
\]
\end{theorem}

\begin{proof}
Direct computation using the FRW Ricci tensor formulae with $k = 1$, $n - 1$ spatial
dimensions, and the cascade identities $\dot{a}^2 + k = 1/a^2$, $\ddot{a}/a = -1/a^4$: $R_{tt} = -(n - 1)\ddot{a}/a =
(n - 1)/a^4$; $R_{ij} = [(n - 2)(\dot{a}^2 + k)/a^2 + \ddot{a}/a]\,g_{ij} = [(n - 2)/a^4 - 1/a^4]\,g_{ij} = (n - 3)/a^4 \cdot g_{ij}$. The
Ricci scalar: $R^{(n)} = g^{tt}R_{tt} + g^{ij}R_{ij} = -(n - 1)/a^4 + (n - 1)(n - 3)/a^4 = (n - 1)(n - 4)/a^4$.
For the normal-normal Einstein component at the equatorial embedding $\psi = \pi/2$ in
$S^{n-1} \subset S^n$: $N_\psi = (1 - t^2)^{-1/2}$, giving $R^{(n)}_{NN} = N_\psi N_\psi R^{(n)}_{\psi\psi} = (n - 3)/a^4$, and $G^{(n)}_{NN} =
R^{(n)}_{NN} - \tfrac{1}{2}g_{NN}R^{(n)} = (n - 3)/a^4 - (n - 1)(n - 4)/(2a^4) = -(n - 2)(n - 5)/(2a^4)$.
\end{proof}

\begin{theorem}[Totally geodesic embedding]
The 4D Lorentzian cascade spacetime
$M^4 = \mathbb{R} \times S^3$ embeds in $M^5 = \mathbb{R} \times S^4$ as a totally geodesic submanifold: $K_{\mu\nu} = 0$. The
same holds at every equatorial embedding in the chain $S^3 \hookrightarrow S^4 \hookrightarrow \cdots \hookrightarrow S^{216}$.
\end{theorem}

\begin{proof}
Parametrise $S^4 = \{(\psi, \theta_i)\}$ with $d\sigma_4^2 = d\psi^2 + \sin^2\psi\,d\sigma_3^2$. The 4D submanifold
sits at $\psi = \pi/2$. The spacelike unit normal is $N = (1 - t^2)^{-1/2}\partial_\psi$. For the temporal
component: $g_{tt} = -1$ is independent of $\psi$, so $K_{tt} = 0$. For the spatial components:
$\partial_\psi g_{\theta_i\theta_j} = 2(1 - t^2)\sin\psi\cos\psi\,\bar{g}_{ij}$; at $\psi = \pi/2$: $\cos(\pi/2) = 0$, so $K_{ij} = 0$. Therefore
$K_{\mu\nu} = 0$ identically. The argument applies at every equatorial embedding by identical
reasoning.
\end{proof}

\begin{corollary}[Hamiltonian constraint]
The Hamiltonian constraint ${}^{(4)}R = -2G^{(5)}_{NN}$
is satisfied identically: $0 = 0$ at every cascade level. The cascade embedding chain is
geometrically self-consistent.
\end{corollary}

\begin{proof}
The scalar Gauss equation for a hypersurface with $K_{\mu\nu} = 0$ gives ${}^{(4)}R = {}^{(5)}R - 2R^{(5)}_{NN}$.
Since $R^{(5)}_{NN} = G^{(5)}_{NN} + \tfrac{1}{2}R^{(5)}$, substituting: ${}^{(4)}R = R^{(5)} - 2(G^{(5)}_{NN} + \tfrac{1}{2}R^{(5)}) = -2G^{(5)}_{NN}$. From
Theorem 14.5 at $n = 5$: $G^{(5)}_{NN} = -(5 - 2)(5 - 5)/(2a^4) = 0$. From Corollary 9.4 at $n = 4$:
$R^{(4)} = 0$. Therefore $0 = -2 \cdot 0 = 0$. \checkmark
\end{proof}

\begin{theorem}[Effective equation of state of the cascade metric]
The 4D Lorentzian
cascade metric has effective equation of state $w = 1/3$ (radiation) in cascade time, at every
cascade level.
\end{theorem}

\begin{proof}
From Theorem 14.5 at $n = 4$ (using $R^{(4)} = 0$, so $G^{(4)}_{\mu\nu} = R^{(4)}_{\mu\nu}$): $G^{(4)}_{tt} = R^{(4)}_{tt} = 3/a^4$;
$G^{(4)}_{ij} = R^{(4)}_{ij} = (4 - 3)/a^4 \cdot g_{ij} = g_{ij}/a^4$. Identifying $\rho_{\text{eff}} \propto G^{(4)}_{tt} = 3/a^4$ and $p_{\text{eff}} \propto G^{(4)}_{ij}/g_{ij} =
1/a^4$, the effective equation of state is:
\[
w_{\text{eff}} = \frac{p_{\text{eff}}}{\rho_{\text{eff}}} = \frac{1/a^4}{3/a^4} = \frac{1}{3}.
\]
\end{proof}

\begin{remark}[Consequences for the matter fraction]
Theorems 14.6 and 14.8 establish
that the geometric Gauss--Codazzi route to the cascade matter fraction $\Omega_m$ is definitively
closed. Three independent exact results confirm this:
\begin{enumerate}
\item $K_{\mu\nu} = 0$ (Theorem 14.6) eliminates the standard extrinsic curvature mechanism for
KK matter generation. The cascade slicing geometry does not source matter via this
channel.
\item $w = 1/3$ (Theorem 14.8) shows the cascade metric contributes radiation, not matter
($w = 0$) or dark energy ($w = -1$), in cascade time. No splitting of the cascade
integrand into core and tail produces the combination ($w = -1$, $w = 0$) consistent
with the Friedmann equation.
\item The Hamiltonian constraint telescopes identically (Corollary 14.7): ${}^{(4)}R = -2G^{(5)}_{NN} =
0 = 0$ at every level, providing geometric consistency without yielding matter content.
\end{enumerate}
The cascade matter fraction $\Omega_m$ is determined by the Bott partition of cascade layers---the sphere-area fraction of non-trivial-phase layers ($d \bmod 8 \in \{5, 6\}$). This requires
three inputs: the propagator phase classification ([2], Corollary 6.5; period-8 refinement in this paper, Section 9), which establishes the
topological partition; the fermion generation structure of the Bott partition ([3]), which
identifies non-trivial-phase layers with matter via the hairy ball zeros and observed fermion
generations; and the uniqueness of sphere areas as cascade content (Corollary 4.8a of [1]),
which forces the energy assignment to be proportional to $\Omega_{d-1}$ with a universal constant.
The matter fraction is a topological rather than geometric quantity; it does not require
the cascade metric and is unaffected by the Gauss--Codazzi results above. The explicit
derivation is given in the cosmological companion.
\end{remark}

\subsection{The information paradox}

The standard black hole information paradox asks how information is preserved when it
crosses a horizon. In the cascade:

\begin{enumerate}
\item \textbf{No sharp horizon.} The compactification at each step is asymptotic: $R_{\text{eff}} = 1/\sqrt{d + 3} >
0$ for all finite $d$. The weight function $(1 - x^2)^{d/2}$ has support everywhere on $(-1, 1)$.
No direction is fully removed; no information is lost.
\item \textbf{Unitary dynamics.} The discrete propagator ([2], Theorem 7.3) preserves the phase
structure. Each slicing step is a unitary operation on the Hilbert space.
\item \textbf{Apparent irreversibility is asymptotic.} From above ($d + 1$ dimensions), a slicing step
looks irreversible: the direction's contribution is suppressed to $O(1/\sqrt{d})$. From below
($d$ dimensions), the boundary carries the full content (boundary dominance). Both
descriptions are correct; they describe the same asymptotic process from different
sides.
\end{enumerate}

There is no paradox because there was never a sharp boundary to lose information
across---only asymptotic compactification. No firewall, no information loss, no complementarity principle needed.

\section{What This Paper Does Not Do}

\begin{itemize}
\item It does not derive Newton's constant $G$ or the local solutions of GR.
\item It does not derive the Standard Model gauge group $\text{SU}(3) \times \text{SU}(2) \times \text{U}(1)$.
\item It does not unify QM and GR in the sense of a quantum gravity theory.
\item It requires one physical identification hypothesis beyond pure mathematics (Definition 2.1): that the cascade geometry is indistinguishable from our universe. Without
it the series is pure mathematics.
\item It uses no result from physics: all remaining inputs are the orthogonality axiom and
classical mathematical theorems (Lovelock, Clifford--Bott, Hurwitz).
\end{itemize}

\section{Open Questions}

\begin{enumerate}
\item \textbf{The Standard Model.} The tower $\{7, 12, 198\}$ (dimensions below $d_0$, between $d_0$
and $d_1$, between $d_1$ and $d_2$) provides natural sub-bundles. The number 12 matches
the generator count of $\text{SU}(3) \times \text{SU}(2) \times \text{U}(1)$.
\item \textbf{Quantum gravity.} The cascade derives QM and GR separately. Whether they
are compatible when both are projected onto 4D simultaneously, and whether the
cascade's discrete dimension structure plays the role of quantised geometry, remains
open.
\item \textbf{Black holes.} Whether the cascade's foliation structure admits horizon formation,
and whether Bekenstein--Hawking entropy $S = A/(4l_{\text{Pl}}^2)$ can be derived from the
sphere-area sequence, is open. The cascade is a theory of sphere areas; the structural
hint is direct.
\item \textbf{Newton's constant.} The cascade produces dimensionless quantities; $G$ enters as
a unit-matching constant from the embedding $M^4 \subset M^{217}$. Deriving $G$ from the
cascade's KK mass spectrum (with compactification radii $R_{\text{eff}}(d) = 1/\sqrt{d + 3}$ now
determined) is the next natural step.
\item \textbf{The cascade Friedmann equation.} The cascade metric (Section 4) foliates a
217-dimensional geometry onto a 4D FRW slice. The projected Friedmann equation
should follow from the Gauss--Codazzi equations applied to the $M^4 \subset M^{217}$ embedding. The Lorentzian Gauss--Codazzi analysis (Section 14.4) establishes that the
cascade metric contributes radiation ($w = 1/3$) and that $K_{\mu\nu} = 0$ eliminates the
extrinsic curvature mechanism. The matter and vacuum energy content of the 4D
Friedmann equation enters via the Bott partition of cascade layers (Remark 14.9), not
via the cascade metric geometry. A complete derivation of the Friedmann equation
connecting the cascade's KK structure to the physical expansion history remains
open.
\item \textbf{The Hubble tension.} The cascade's parameter triple $(H_0, \Omega_m, \Omega_b)$ gives a sound
horizon smaller than the Planck-inferred value. The cascade's $H_0$ sits between
the two principal measurements. The full cosmological derivation is a separate
programme.
\end{enumerate}

\begin{thebibliography}{12}

\bibitem{paper1} {[Author]}, \textit{Scale Variance from Orthogonality: How the Unit Ball Generates $10^{120}$ Orders of Magnitude}, March 2026.

\bibitem{paper2} {[Author]}, \textit{Quantum Mechanics from the Cascade: Effective Theory of a 4-Dimensional Observer in the Sphere-Area Geometry}, March 2026.

\bibitem{paper4a} {[Author]}, \textit{The Standard Model from the Cascade: Part IVa---Gauge Group, Symmetry Breaking, and Three Generations from Bott Periodicity and Hairy Ball Zeros}, March 2026.

\bibitem{lovelock1} D.~Lovelock, ``The Einstein tensor and its generalisations,'' \textit{Journal of Mathematical Physics} \textbf{12}, 498--501 (1971).

\bibitem{lovelock2} D.~Lovelock, ``The four-dimensionality of space and the Einstein tensor,'' \textit{Journal of Mathematical Physics} \textbf{13}, 874--876 (1972).

\bibitem{lounesto} P.~Lounesto, \textit{Clifford Algebras and Spinors}, 2nd ed., Cambridge University Press, 2001.

\bibitem{hurwitz} A.~Hurwitz, ``\"Uber die Komposition der quadratischen Formen,'' \textit{Mathematische Annalen} \textbf{88}, 1--25 (1923).

\bibitem{kaluza} T.~Kaluza, ``Zum Unit\"atsproblem der Physik,'' \textit{Sitzungsberichte der Preussischen Akademie der Wissenschaften}, 966--972 (1921).

\bibitem{klein} O.~Klein, ``Quantentheorie und f\"unfdimensionale Relativit\"atstheorie,'' \textit{Zeitschrift f\"ur Physik} \textbf{37}, 895--906 (1926).

\bibitem{padmanabhan} T.~Padmanabhan, ``Some aspects of the cosmological constant problem,'' \textit{Physics Reports} \textbf{380}, 235--320 (2003).

\bibitem{weinberg} S.~Weinberg, ``The cosmological constant problem,'' \textit{Reviews of Modern Physics} \textbf{61}, 1--23 (1989).

\bibitem{adm} R.~Arnowitt, S.~Deser, and C.~W.~Misner, ``The dynamics of general relativity,'' in \textit{Gravitation: An Introduction to Current Research}, Wiley, 1962.

\end{thebibliography}

\end{document}
